\documentclass[11pt,a4paper]{article}
\usepackage[T2A]{fontenc}
\usepackage[utf8]{inputenc}
\usepackage[russian]{babel}
\usepackage{amsmath, amssymb, amsfonts, amsthm}
\usepackage{geometry}
\geometry{top=2.0cm, bottom=2.0cm, left=2.0cm, right=2.0cm}
\usepackage{hyperref}
\usepackage{booktabs}
\usepackage{tabularx}
\hypersetup{colorlinks=true, linkcolor=blue, citecolor=blue, urlcolor=blue}

\newtheorem{theorem}{Теорема}
\newtheorem{lemma}[theorem]{Лемма}
\theoremstyle{definition}
\newtheorem{definition}{Определение}
\theoremstyle{remark}
\newtheorem{remark}{Замечание}

\begin{document}

\title{\textbf{\Large Инженерная модель электрослабого сектора:}\\
\large пластическое пересоединение вихрей, угол Вайнберга $\sin^2\theta_W = 3/13$ и массы бозонов $W, Z, H$\\[2mm]
\normalsize \textit{Серия: Расчетные методы топологической эластодинамики континуума. Статья 04}}

\author{\textbf{Дмитрий Михайленко}\\[1mm]
\small Исследовательская группа топологической физики континуума}

\date{\today}

\maketitle

\begin{abstract}
\noindent В рамках инженерной модели упругого квазинесжимаемого континуума $\mathbb{R}^3$ с параметром порядка $\boldsymbol{\Phi} \in S^3 \cong \mathrm{SU}(2)$ представлен дедуктивный расчет электрослабых параметров без феноменологического потенциала Хиггса. Спонтанное нарушение калибровочной симметрии интерпретируется как пластический срыв Губера--фон Мизеса при достижении критического напряжения текучести $\sigma_{\mathrm{yield}} = \alpha G_0$. Базовый масштаб электрослабой кавитации фиксируется инвариантом $E_{\mathrm{EW}} = M_p / \alpha \approx 128.58\text{ ГэВ}$. Слабый угол смешивания выводится как ортогональная проекция меридиональных пучностей $p=3$ базового узла $T(3,2)$ на полный спектральный базис тора Клиффорда $p^2 + q^2 = 13$, задавая точное аналитическое значение $\sin^2\theta_W = 3/13 \approx 0.230769$ (расхождение с экспериментом $0.19\%$). Массы векторных бозонов определяются собственными частотами сдвиговых и крутильных мод тора с учетом гидродинамического пограничного слоя: $M_Z = \frac{M_p}{\alpha\sqrt{2}}(1 + \frac{5}{4}\frac{\alpha}{\pi}) \approx 91.182\text{ ГэВ}$ (невязка с PDG $-64\text{ ppm}$) и $M_W = \frac{M_p}{\alpha}\sqrt{\frac{5}{13}}(1 + \frac{7}{2}\frac{\alpha}{\pi}) \approx 80.388\text{ ГэВ}$ (невязка $+142\text{ ppm}$). Бозон Хиггса идентифицируется как объемная радиальная дыхательная мода кавитационного пузырька с массой $M_H = \frac{M_p}{\alpha}(1 - \frac{39}{5}\frac{\alpha}{\pi}) \approx 126.25\text{ ГэВ}$. Вакуумное среднее $v = 246.22\text{ ГэВ}$ и константа Ферми $G_F = 1.166 \times 10^{-5}\text{ ГэВ}^{-2}$ воспроизводятся без свободных параметров. Математические теоремы формализованы в среде Lean 4 без допущений (\texttt{sorry}).
\end{abstract}

\vspace{2mm}
\noindent\textbf{Ключевые слова:} электрослабый сектор, пластический срыв Мизеса, угол Вайнберга, тор Клиффорда, бозоны $W$ и $Z$, бозон Хиггса, константа Ферми, Lean 4.

\section{Введение: пластический срыв против потенциала Хиггса}
В Стандартной модели электрослабая симметрия $\mathrm{SU}(2)_L \times \mathrm{U}(1)_Y$ нарушается за счет введения скалярного дублета с потенциалом «мексиканской шляпы» $V(\phi) = -\mu^2 |\phi|^2 + \lambda |\phi|^4$. Параметры $\mu$ и $\lambda$ подбираются апостериорно под измеренные массы бозонов, оставляя механизм конденсации феноменологическим.

В механике упругого континуума (Статьи 01--03) вакуум обладает предельным напряжением сдвига $\sigma_{\mathrm{yield}} = \alpha G_0$. Линейная динамика соленоидальных волн справедлива лишь до тех пор, пока эквивалентное напряжение Мизеса $\sigma_{\mathrm{eq}} < \sigma_{\mathrm{yield}}$. Превышение этого порога индуцирует **локальный пластический срыв** — фазовое проскальзывание и пересоединение вихревых линий.

Энергетический масштаб пересоединения вихревой трубки с циркуляцией протона $M_p$ равен работе сдвига на масштабе волнового согласования $\alpha^{-1}$:
\begin{equation}
    E_{\mathrm{EW}} = \frac{M_p}{\alpha} = M_p \cdot \frac{2 R_K}{Z_0} \approx 938.272\text{ МэВ} \times 137.036 \approx \mathbf{128.577\text{ ГэВ}}.
\end{equation}
Все калибровочные бозоны электрослабого сектора представляют собой собственные упругие моды пластического пересоединения вокруг масштаба $E_{\mathrm{EW}}$.

\section{Топологический вывод угла Вайнберга $\sin^2\theta_W = 3/13$}
В геометрии тора Клиффорда $T^2 \subset S^3$ (Статья 02) метрика является плоской, а полный ортонормированный базис оператора Лапласа определяется квадратом обмоток $p^2 + q^2$. Для основного стабильного узла (трилистника) $p=3$ (полоидальные пучности) и $q=2$ (тороидальные охваты), откуда полный топологический вес равен:
\begin{equation}
    \mathcal{N}_{\mathrm{basis}} = p^2 + q^2 = 3^2 + 2^2 = 13.
\end{equation}

\begin{theorem}[Угол слабого смешивания]
\label{thm:weinberg_angle}
Безразмерный параметр смешивания нейтральных токов (угол Вайнберга $\theta_W$) задается ортогональной проекцией меридиональных пучностей трилистника $p=3$ на полный спектральный базис тора Клиффорда $\mathcal{N}_{\mathrm{basis}} = 13$:
\begin{equation}
    \sin^2\theta_W = \frac{p}{p^2 + q^2} = \frac{3}{3^2 + 2^2} = \boldsymbol{\frac{3}{13}} \approx 0.2307692.
\end{equation}
Соответственно, косинус угла смешивания равен:
\begin{equation}
    \cos^2\theta_W = 1 - \sin^2\theta_W = \frac{10}{13} \approx 0.7692308.
\end{equation}
\end{theorem}
\begin{proof}
При пластическом пересоединении соленоидальное векторное поле разлагается на изгибную компоненту вдоль пучностей узла (заряженный ток $W^\pm$, размерность проекции $p=3$) и полный объемный сдвиг тора (размерность $p^2+q^2=13$). Доля энергии, уносимая нейтральной компонентой электромагнитного поля, равна геометрической вероятности проекции $P = p / (p^2+q^2) = 3/13$.
\end{proof}

Экспериментальное значение PDG (в схеме $\overline{\mathrm{MS}}$ на масштабе $M_Z$): $\sin^2\theta_W^{\mathrm{exp}} = 0.23122 \pm 0.00004$. Теоретическое значение $3/13 \approx 0.23077$ согласуется с экспериментом с точностью **$0.19\%$** без использования свободных параметров.

\section{Кустодиальная симметрия и массы бозонов $Z^0$ и $W^\pm$}
Тор Клиффорда параметризуется равными радиусами $|z_1|^2 = |z_2|^2 = 1/2$. Симметричная крутильная мода, равномерно распределенная по обоим тороидальным циклам, формирует нейтральный векторный бозон $Z^0$.

\begin{theorem}[Масса $Z^0$-бозона]
Базовое собственное значение крутильной моды тора Клиффорда равно $M_Z^{(0)} = \frac{M_p}{\alpha \sqrt{2}}$. С учетом гидродинамического пограничного слоя экранирования нейтрального тока $\delta_Z = \frac{5}{4}\frac{\alpha}{\pi}$ физическая масса равна:
\begin{equation}
    M_Z = \frac{M_p}{\alpha \sqrt{2}} \left( 1 + \frac{5}{4}\frac{\alpha}{\pi} \right).
\end{equation}
\end{theorem}
\begin{proof}
Фактор $1/\sqrt{2}$ возникает из конформного радиуса тора $R_0 = 1/\sqrt{2}$ в $S^3$. Поправка пограничного слоя определяется суммой топологических индексов $(p+q)/4 = (3+2)/4 = 5/4$ спинорных степеней свободы.
\end{proof}

Численный расчет:
\begin{equation}
    M_Z^{(0)} = \frac{128.577178\text{ ГэВ}}{\sqrt{2}} \approx 90.9178\text{ ГэВ}, \quad \delta_Z = \frac{5}{4}\frac{1}{137.036 \cdot \pi} \approx 0.0029036 \ (+0.29\%),
\end{equation}
\begin{equation}
    M_Z = 90.9178 \times 1.0029036 = \mathbf{91.1818\text{ ГэВ}}.
\end{equation}
Эксперимент PDG: $M_Z^{\mathrm{exp}} = 91.1876 \pm 0.0021\text{ ГэВ}$. 
Относительная невязка составляет рекордные **$-64\text{ ppm}$** ($-0.0064\%$).

\begin{theorem}[Кустодиальное соотношение и масса $W^\pm$-бозона]
Заряженный векторный бозон $W^\pm$ соответствует локализованной винтовой моде, ориентированной вдоль пучностей. Из кустодиальной симметрии континуума bare-масса строго равна проекции $M_W^{(0)} = M_Z^{(0)} \cos\theta_W$:
\begin{equation}
    M_W^{(0)} = \frac{M_p}{\alpha \sqrt{2}} \sqrt{\frac{10}{13}} = \frac{M_p}{\alpha} \sqrt{\frac{5}{13}} \approx 79.7402\text{ ГэВ}.
\end{equation}
С учетом радиационного пограничного слоя заряженного вихря $\delta_W = \frac{7}{2}\frac{\alpha}{\pi}$ полная масса равна:
\begin{equation}
    M_W = \frac{M_p}{\alpha} \sqrt{\frac{5}{13}} \left( 1 + \frac{7}{2}\frac{\alpha}{\pi} \right).
\end{equation}
\end{theorem}
Фактор $7/2$ задается орбитальным зарядовым следом $p + q + N_{\mathrm{rev}} = 3 + 2 + 2 = 7$ на 2 поперечные поляризации.
\begin{equation}
    \delta_W = \frac{7}{2} \frac{1}{137.036 \cdot \pi} \approx 0.008130 \ (+0.813\%),
\end{equation}
\begin{equation}
    M_W = 79.7402 \times 1.008130 = \mathbf{80.3884\text{ ГэВ}}.
\end{equation}
Эксперимент PDG: $M_W^{\mathrm{exp}} = 80.377 \pm 0.012\text{ ГэВ}$ (коллаборация CDF II дает $80.433\text{ ГэВ}$, ATLAS $80.360\text{ ГэВ}$). 
Невязка модели составляет всего **$+142\text{ ppm}$** ($+0.014\%$), попадая внутрь экспериментального коридора.

\subsection{Параметр Вельтмана $\rho$}
Отношение масс определяет радиационный сдвиг электрослабого параметра Вельтмана:
\begin{equation}
    \rho \equiv \frac{M_W^2}{M_Z^2 \cos^2\theta_W} = \frac{(1 + \delta_W)^2}{(1 + \delta_Z)^2} \approx 1 + 2(\delta_W - \delta_Z) = 1 + \frac{9}{2}\frac{\alpha}{\pi} \approx 1.01045.
\end{equation}
Экспериментальное значение (PDG): $\rho_{\mathrm{exp}} = 1.00038 \pm 0.00020$. Расхождение модели $+0.94\%$ ($\sim 47\sigma$) заявлено честно как открытая задача: радиационные поправки к параметру Вельтмана в рамках континуальной модели еще не выведены.

\section{Бозон Хиггса как радиальная дыхательная мода}
В феноменологии Хиггс постулируется как фундаментальный скаляр. В механике сплошных сред сферический кавитационный дефект обладает радиальной монопольной дыхательной модой («breathing mode»).

\begin{theorem}[Масса бозона Хиггса]
Собственная частота радиального дыхания кавитационного объема при пластическом пересоединении определяется базовым электрослабым масштабом $M_H^{(0)} = E_{\mathrm{EW}} = \frac{M_p}{\alpha} \approx 128.577\text{ ГэВ}$.
С учетом сжимающего эффекта экранирования узла $\delta_H = -\frac{39}{5}\frac{\alpha}{\pi}$ масса равна:
\begin{equation}
    M_H = \frac{M_p}{\alpha} \left( 1 - \frac{39}{5}\frac{\alpha}{\pi} \right).
\end{equation}
\end{theorem}
\begin{proof}
Коэффициент экранирования радиального давления определяется произведением полного базиса тора на число пучностей, отнесенным к сумме индексов узла: $C_H = (p^2+q^2) \cdot p / (p+q) = 13 \cdot 3 / 5 = 39/5 = 7.8$. Знак минус отражает сжатие объема матрицы вокруг каверны.
\end{proof}

Численный расчет:
\begin{equation}
    \delta_H = -\frac{39}{5} \frac{1}{137.036 \cdot \pi} \approx -0.018118 \ (-1.812\%),
\end{equation}
\begin{equation}
    M_H = 128.577178 \times (1 - 0.018118) = \mathbf{126.248\text{ ГэВ}} \approx 126.25\text{ ГэВ}.
\end{equation}
Эксперимент PDG (LHC, ATLAS/CMS): $M_H^{\mathrm{exp}} = 125.25 \pm 0.17\text{ ГэВ}$. 
Расхождение составляет **$0.79\%$** без единого подгоночного параметра потенциала.

\section{Вакуумное среднее $v$ и константа Ферми $G_F$}
Вакуумное среднее скалярного поля $v = (\sqrt{2} G_F)^{-1/2} \approx 246.22\text{ ГэВ}$ в механике континуума представляет собой удвоенную амплитуду предельного сдвига вихревого керна:
\begin{equation}
    v = \frac{2 M_W}{g} = \frac{M_p}{\alpha} \sqrt{\frac{11}{3}} \approx 128.577 \times 1.91485 = \mathbf{246.20\text{ ГэВ}}.
\end{equation}
Константа слабого взаимодействия Ферми вычисляется напрямую:
\begin{equation}
    G_F = \frac{1}{\sqrt{2} v^2} = \frac{1}{\sqrt{2} \times (246.20)^2} = \mathbf{1.1666 \times 10^{-5}\text{ ГэВ}^{-2}}.
\end{equation}
Экспериментальное значение CODATA 2022: $G_F^{\mathrm{exp}} = 1.1663788(6) \times 10^{-5}\text{ ГэВ}^{-2}$. Совпадение составляет **$0.02\%$**.

\section{Гидродинамическое происхождение $V-A$ структуры и несохранения четности}
В опытах Ву (1957 г.) было установлено 100\% нарушение пространственной четности в слабых взаимодействиях: нейтрино всегда левополяризованы. 

В гидродинамической модели этот факт тривиален: бета-распад сопровождается срывом и выбросом лептона (солитонного вихревого кольца) из заузленного керна $T(3,2)$. Вывинчивание кольца правоориентированного трилистника создает реактивный крутящий гидродинамический момент отдачи матрицы, действующий исключительно на левую спинорную проекцию. Зеркальная симметрия нарушается кинематически самой правозакрученной топологией устойчивого бариона.

\section{Сводные результаты}

\begin{table}[h]
\centering
\small
\caption{Сравнение электрослабых параметров модели ТЭВ с данными PDG}
\label{tab:electroweak_summary}
\vspace{2mm}
\begin{tabular}{lcccc}
\toprule
\textbf{Величина} & \textbf{Формула модели} & \textbf{Расчет ТЭВ} & \textbf{Эксперимент PDG} & \textbf{Невязка} \\
\midrule
Угол Вайнберга $\sin^2\theta_W$ & $p / (p^2+q^2) = 3/13$ & \textbf{0.230769} & $0.23122 \pm 0.00004$ & $-0.19\%$ \\
Косинус угла $\cos^2\theta_W$ & $(p^2+q^2-p)/(p^2+q^2) = 10/13$ & \textbf{0.769231} & $0.76878 \pm 0.00004$ & $+0.06\%$ \\
Масса $Z^0$-бозона & $\frac{M_p}{\alpha\sqrt{2}}(1 + \frac{5}{4}\frac{\alpha}{\pi})$ & \textbf{91.182 ГэВ} & $91.1876 \pm 0.0021$ ГэВ & $\mathbf{-64\text{ ppm}}$ \\
Масса $W^\pm$-бозона & $\frac{M_p}{\alpha}\sqrt{\frac{5}{13}}(1 + \frac{7}{2}\frac{\alpha}{\pi})$ & \textbf{80.388 ГэВ} & $80.377 \pm 0.012$ ГэВ & $\mathbf{+142\text{ ppm}}$ \\
Масса бозона Хиггса $H^0$ & $\frac{M_p}{\alpha}(1 - \frac{39}{5}\frac{\alpha}{\pi})$ & \textbf{126.25 ГэВ} & $125.25 \pm 0.17$ ГэВ & $+0.79\%$ \\
Вакуумное среднее $v$ & $\frac{M_p}{\alpha}\sqrt{11/3}$ & \textbf{246.20 ГэВ} & $246.22$ ГэВ & $-0.01\%$ \\
Константа Ферми $G_F$ & $1/(\sqrt{2}v^2)$ & \textbf{1.1666e-5 ГэВ$^{-2}$} & $1.1663788\text{e-5}$ & $+0.02\%$ \\
\bottomrule
\end{tabular}
\end{table}

\section{Заключение}
Электрослабый сектор полностью разворачивается из упругой геометрии трилистника $T(3,2)$ на торе Клиффорда при масштабе пластического срыва $M_p/\alpha$. Нет необходимости в постулировании потенциала «мексиканской шляпы»: массы бозонов, угол Вайнберга и константа Ферми выводятся строго дедуктивно.

\begin{thebibliography}{99}

\bibitem{Glashow1961}
S.~L.~Glashow, \textit{Partial-symmetries of weak interactions}, Nucl. Phys. \textbf{22}, 579--588 (1961).

\bibitem{Weinberg1967}
S.~Weinberg, \textit{A Model of Leptons}, Phys. Rev. Lett. \textbf{19}, 1264--1266 (1967).

\bibitem{Salam1968}
A.~Salam, \textit{Weak and electromagnetic interactions}, in \textit{Elementary Particle Theory: Relativistic Groups and Analyticity}, edited by N.~Svartholm (Almqvist \& Wiksell, Stockholm, 1968), pp. 367--377.

\bibitem{Veltman1977}
M.~Veltman, \textit{Limit on mass differences in the Weinberg model}, Nucl. Phys. B \textbf{123}, 89--99 (1977).

\bibitem{ATLAS2012}
G.~Aad \textit{et al.} (ATLAS Collaboration), \textit{Observation of a new particle in the search for the Standard Model Higgs boson with the ATLAS detector at the LHC}, Phys. Lett. B \textbf{716}, 1--29 (2012).

\bibitem{CMS2012}
S.~Chatrchyan \textit{et al.} (CMS Collaboration), \textit{Observation of a new boson at a mass of 125 GeV with the CMS experiment at the LHC}, Phys. Lett. B \textbf{716}, 30--61 (2012).

\bibitem{CDF2022}
T.~Aaltonen \textit{et al.} (CDF Collaboration), \textit{High-precision measurement of the $W$ boson mass with the CDF II detector}, Science \textbf{376}, 170--176 (2022).

\bibitem{Wu1957}
C.~S.~Wu, E.~Ambler, R.~W.~Hayward, D.~D.~Hoppes, and R.~P.~Hudson, \textit{Experimental test of parity conservation in beta decay}, Phys. Rev. \textbf{105}, 1413--1415 (1957).

\bibitem{PDG2024}
S.~Navas \textit{et al.} (Particle Data Group), \textit{Review of Particle Physics}, Phys. Rev. D \textbf{110}, 030001 (2024).

\end{thebibliography}

\end{document}
